\section{引言}\label{sec:intro}
LLM Agent能够辩论、批评、投票并综合彼此输出，但聚合结果的价值取决于产生它的行为。短答案、独立推理报告、追加验证和模仿同伴具有不同成本与信息价值。多数评测只观察最终答案，使系统无法区分真实审议与廉价一致。

这首先是机制设计问题，其次才是聚合问题。设计者必须在Agent行动前宣布何种行为获得信用，Agent再决定预期信用是否值得资源投入。对既有回答做回溯加权不能建立追随者响应，因为生成策略未观察该规则；改变拓扑也只能控制信息流向，不能决定Agent是否生产可独立检查的信息。

本文提出\textbf{博弈驱动的可观测语义共识机制}，并以Assumption、Evidence、Inference和Conclusion四个可解析字段构成的AEIC报告作为观测接口。求解前，追随者观察包含信用规则、行动成本、信息策略和验证预算的公开合同，并从Low、AEIC、Verify和Imitate中选择。行动选择真实生成路径；报告、资源、信用、成本和信誉更新均被记录。这把“先承诺、后响应”的Stackelberg直觉转化为可证伪实验。AEIC负责使过程可解析和可结算，机制处理则是行动前合同及其与真实资源路径的绑定。

确认性实验在三个任务域的450道留出题上配对比较三种规则。Static AEIC将高资源行动率从0.46\%提高到56.93\%，联合执行配置下的准确率从68.00\%提高到82.22\%；Outcome Only仅提高到3.08\%，准确率无显著变化。因此现有结果支持局部可读过程规则改变行动策略，但由于条件同时改变执行提示、Token和后端路由，质量增益不能单独归因于AEIC格式或信用文本。

本文还用奖励--成本扫描、四类信息环境、信号与历史消融、额外Seed、不同面板规模和留出GLM后端检验解释。GLM未响应同一合同，这一负结果排除普适性强主张，并将后端合同校准确定为开放问题。

本文贡献为：\begin{enumerate}
\item 可执行Leader--Follower协议：规则先于行动，行动控制真实资源路径；
\item 配对证据表明基于AEIC观测接口的过程合同相对仅结果奖励能够改变资源行动，并在联合执行配置下伴随集体准确率改善；
\item 完整冻结的八组评测，覆盖剂量响应、信息控制、消融、稳健性和明确的不迁移边界。
\end{enumerate}
